\subsection{Milestone 3 -- Risultati dell'analisi}

\subsubsection{Dataset e predizioni}

La Tabella~\ref{tab:r3-datasets} riporta il numero di istanze e i bug predetti per ciascun
dataset. Le predizioni superano sempre gli actual (A: $+6$, B+: $+6$, B: $+18$; C è esatto):
l'effetto è coerente con SMOTE, che sposta la soglia decisionale verso ``buggy'' producendo
una leggera sovrastima. Questo non implica che le istanze predette coincidano
con quelle realmente buggy: il modello può classificare erroneamente classi clean come buggy
e mancarne altre effettivamente difettose. BClassifierA è addestrato su A, quindi le predizioni su A sono dentro il campione
di training: il modello ha già visto quelle istanze e i risultati sono gonfiati.
Invece per la M2 (AUC $0{,}94$, Kappa $0{,}61$) ogni
istanza è stata predetta da un modello che non l'aveva mai vista in training (cross-validation):
i due valori non sono confrontabili. Il bias introdotto dall'aver visto A agisce in modo
identico su B+ e B, cancellandosi nel confronto dell'analisi what-if.

\subsubsection{Metriche di prevenzione}

I risultati sono riportati in Tabella~\ref{tab:r3-prevention}.

Su 1.125 classi predette buggy in B+, solo \textbf{63} (5,60\%) smetterebbero di essere
predette buggy se i code smell fossero zero. Un dato inatteso: B predice \textbf{1.137}
bug contro le 1.125 di B+ ($+12$ netto). L'azzeramento degli smell genera 63 istanze
prevenute (BUGGY $\to$ CLEAN) ma anche 75 istanze nella direzione opposta (CLEAN $\to$
BUGGY).

\subsubsection{Confronto con il paper di riferimento}

Il paper "What if I had no smells" presentato a lezione riporta
una riduzione del \textbf{42\%} dei file difettosi predetti passando da B+ a B, contro
il nostro 5,60\%. La metodologia è identica (A/B+/B/C, stesso criterio, modello addestrato
su A); le differenze risiedono nel contesto sperimentale:

\begin{enumerate}

\item \textbf{Strumento e ruleset.} Il paper usa \textbf{SonarQube} , in questo progetto si usa invece
\textbf{PMD}, dove alcune "regole" rilevano violazioni di stile puro
(naming, import superflui...) che crescono con la
dimensione del file. Questo introduce una correlazione artificiale tra NSmells e LOC, che
si propaga anche alle metriche di processo cumulative (LOC\_Touched, Churn, NR), riducendo
l'indipendenza informativa di NSmells.

\item \textbf{Contributo marginale di NSmells.} Nel paper, NSmells presenta correlazione
quasi nulla con LOC\_touched ($-0{,}01$), NR ($0{,}02$): gli smell sono \textit{indipendenti}
dalle metriche e il loro contributo al modello è aggiuntivo. In questo progetto
(Tabella~\ref{tab:r3-correlation}), NSmells correla $\rho = 0{,}78$ con LOC,
$\rho > 0{,}65$ con Churn, LOC\_Touched, Weighted\_Age.

\end{enumerate}

\subsubsection{Analisi di ridondanza (M3.1)}

\paragraph{Correlazione di Spearman.}
Sette feature su 18 presentano $\rho > 0{,}60$ con NSmells: LOC ($0{,}78$), Max\_Churn
($0{,}71$), LOC\_Touched ($0{,}69$), Churn ($0{,}69$), Weighted\_Age ($0{,}69$),
Max\_LOC\_Added ($0{,}66$), LOC\_Added ($0{,}66$). Le uniche feature con bassa correlazione verso NSmells
sono EXP ($0{,}04$) e SEXP ($0{,}13$), che misurano l'esperienza dell'autore e sono
indipendenti dalla dimensione del file.

\paragraph{Ablation study.}
La Tabella~\ref{tab:r3-ablation} confronta le prestazioni di RF + No FS + SMOTE con e
senza NSmells in 10$\times$10-fold CV su A. Il $\Delta$ è essenzialmente zero su tutte
le metriche.

\paragraph{Sintesi M3.1.}
NSmells correla fortemente con le metriche di processo ($\rho$ fino a $0{,}78$ con LOC)
\textit{e} rimuoverla non degrada il modello ($\Delta \approx 0$). Insieme, queste evidenze
mostrano che, in questo progetto e con queste condizioni, NSmells risulta poco informativo
rispetto alle metriche di processo già presenti nel dataset. Questo non è una conclusione
generale: lo scarso contributo di NSmells dipende dal ruleset PMD scelto, dalla sua
correlazione con LOC in Apache OpenJPA, e dalle caratteristiche specifiche del progetto
analizzato. Su un progetto diverso, con uno strumento diverso o un profilo di sviluppo
diverso, NSmells potrebbe avere un contributo marginale significativamente maggiore ---
come suggerito dal 42\% del paper. Il principale driver di bugginess osservato in
Apache OpenJPA è la storia di modifica del codice, non la qualità strutturale misurata
tramite PMD.
